\documentclass[UTF8,11pt]{ctexart}
\usepackage[a4paper,margin=24mm]{geometry}
\usepackage{amsmath,amssymb,amsthm,mathtools,mathrsfs}
\usepackage[all]{xy}
\usepackage{xurl,hyperref,enumitem}
\usepackage{tikz}
\usetikzlibrary{arrows.meta,cd,decorations.markings,calc,patterns}
\hypersetup{hidelinks,breaklinks=true}
\setlength{\emergencystretch}{3em}
\allowdisplaybreaks
\newtheorem*{theorem}{Theorem}
\newtheorem*{thm}{Theorem}
\newtheorem*{lemma}{Lemma}
\newtheorem*{proposition}{Proposition}
\newtheorem*{prop}{Proposition}
\newtheorem*{corollary}{Corollary}
\newtheorem*{cor}{Corollary}
\newtheorem*{claim}{Claim}
\theoremstyle{definition}
\newtheorem*{definition}{Definition}
\newtheorem*{defn}{Definition}
\newtheorem*{remark}{Remark}
\newtheorem*{example}{Example}
\newcommand{\R}{\mathbb R}
\newcommand{\C}{\mathbb C}
\newcommand{\Z}{\mathbb Z}
\newcommand{\Q}{\mathbb Q}
\newcommand{\N}{\mathbb N}
\newcommand{\F}{\mathbb F}
\newcommand{\D}{\mathbb D}
\newcommand{\bB}{\mathbb B}
\newcommand{\bC}{\mathbb C}
\newcommand{\bR}{\mathbb R}
\newcommand{\bZ}{\mathbb Z}
\newcommand{\bN}{\mathbb N}
\newcommand{\bQ}{\mathbb Q}
\newcommand{\bD}{\mathbb D}
\newcommand{\bF}{\mathbb F}
\newcommand{\CP}{\mathbb{CP}}
\newcommand{\RP}{\mathbb{RP}}
\newcommand{\sO}{\mathscr O}
\newcommand{\sI}{\mathscr I}
\newcommand{\sF}{\mathscr F}
\newcommand{\sG}{\mathscr G}
\newcommand{\sH}{\mathscr H}
\newcommand{\sR}{\mathscr R}
\newcommand{\sM}{\mathscr M}
\newcommand{\sV}{\mathscr V}
\newcommand{\sU}{\mathscr U}
\DeclareMathOperator{\Hom}{Hom}
\DeclareMathOperator{\Ext}{Ext}
\DeclareMathOperator{\Tor}{Tor}
\DeclareMathOperator{\Spec}{Spec}
\DeclareMathOperator{\Proj}{Proj}
\DeclareMathOperator{\supp}{supp}
\DeclareMathOperator{\im}{im}
\DeclareMathOperator{\id}{id}
\DeclareMathOperator{\Tr}{Tr}
\DeclareMathOperator{\tr}{tr}
\DeclareMathOperator{\rank}{rank}
\DeclareMathOperator{\ord}{ord}
\DeclareMathOperator{\dist}{dist}
\DeclareMathOperator{\End}{End}
\DeclareMathOperator{\Aut}{Aut}
\DeclareMathOperator{\Gal}{Gal}
\DeclareMathOperator{\vol}{vol}
\DeclareMathOperator{\Cl}{Cl}
\DeclareMathOperator{\coker}{coker}
\DeclareMathOperator{\codim}{codim}
\DeclareMathOperator{\divergence}{div}
\DeclareMathOperator{\Ric}{Ric}
\DeclareMathOperator{\grad}{grad}
\DeclareMathOperator{\Hess}{Hess}
\newcommand{\abs}[1]{\left\lvert#1\right\rvert}
\newcommand{\sabs}[1]{\lvert#1\rvert}
\newcommand{\babs}[1]{\bigl\lvert#1\bigr\rvert}
\newcommand{\Babs}[1]{\Bigl\lvert#1\Bigr\rvert}
\newcommand{\bbabs}[1]{\biggl\lvert#1\biggr\rvert}
\newcommand{\BBabs}[1]{\Biggl\lvert#1\Biggr\rvert}
\newcommand{\norm}[1]{\left\lVert#1\right\rVert}
\newcommand{\snorm}[1]{\lVert#1\rVert}
\newcommand{\bnorm}[1]{\bigl\lVert#1\bigr\rVert}
\newcommand{\Bnorm}[1]{\Bigl\lVert#1\Bigr\rVert}
\newcommand{\bbnorm}[1]{\biggl\lVert#1\biggr\rVert}
\newcommand{\BBnorm}[1]{\Biggl\lVert#1\Biggr\rVert}
\newcommand{\widebar}[1]{\overline{#1}}
\newcommand{\nicefrac}[2]{\frac{#1}{#2}}
\newcommand{\linnprod}[2]{\langle#1,#2\rangle}
\newcommand{\myindex}[1]{#1}
\newcommand{\glsadd}[1]{}
\newcommand{\avoidbreak}{}
\newcommand{\sourceref}[1]{\textnormal{[原文：\nolinkurl{#1}]}}
\newcommand{\thmref}[1]{\sourceref{#1}}
\newcommand{\lemmaref}[1]{\sourceref{#1}}
\newcommand{\propref}[1]{\sourceref{#1}}
\newcommand{\corref}[1]{\sourceref{#1}}
\newcommand{\defnref}[1]{\sourceref{#1}}
\newcommand{\exampleref}[1]{\sourceref{#1}}
\newcommand{\exerciseref}[1]{\sourceref{#1}}
\newcommand{\figureref}[1]{\sourceref{#1}}
\newcommand{\sectionref}[1]{\sourceref{#1}}
\newcommand{\appendixref}[1]{\sourceref{#1}}
\newcommand{\stacksref}[2]{\href{https://stacks.math.columbia.edu/tag/#1}{#2 [#1]}}


\newcommand{\E}{\mathbb E}
\newcommand{\Prob}{\mathbb P}
\newcommand{\Reff}{\mathcal R_{\mathrm{eff}}}
\newcommand{\eps}{\varepsilon}
\DeclareMathOperator{\diam}{diam}
\DeclareMathOperator{\Var}{Var}
\DeclareMathOperator{\Crit}{Crit}
\newcommand{\dd}{\,\mathrm d}

\begin{document}
\section*{059\quad 单射全纯映射的Jacobian非退化性}
\subsection*{来源与计数目标}
Ji\v{r}\'i Lebl, \emph{Tasty Bits of Several Complex Variables}, Chapter 2，``Riemann extension, zero sets, and injective maps''。\texttt{scv.tex} 中 \texttt{thm:regptsdense} 和 \texttt{thm:injective} 及其全部证明；读取第3810--4125行所覆盖的题证。commit \texttt{1149c5827e3ee1e25cd0939bfa846ca441517e62}，获取日期2026-09-28。
\par\url{https://github.com/jirilebl/scv/blob/1149c5827e3ee1e25cd0939bfa846ca441517e62/scv.tex}
\par\textbf{本文件只计一个问题：}证明同维单射全纯映射的 Jacobian 处处非零，并得到全纯双射的逆全纯；正则零点稠密性只作支撑，不另计题。
\subsection*{作者命题与人类证明}
\begin{theorem}[Regular points of a zero set]\label{regptsdense}
Suppose $U\subset\C^n$ is a domain, $f\in\sO(U)$, and $f$ is not identically zero. Let $N=f^{-1}(0)$. Then there exists an open and dense (subspace topology) subset $N_{\mathit{reg}}\subset N$ such that at each $p\in N_{\mathit{reg}}$, after possibly reordering variables, $N$ can be locally (that is, in some neighborhood) written as
\[z_n=g(z_1,\ldots,z_{n-1})\]
for a holomorphic function $g$.
\end{theorem}
\begin{proof}
If $N$ is locally a graph at $p$, then it is a graph for every point of $N$ near $p$. So $N_{\mathit{reg}}$ is open. If for every point $p_0\in N$ and every neighborhood $W$ of $p_0$, we show that $N\cap W$ has a regular point, then $N_{\mathit{reg}}$ is dense. Replacing $N$ with $N\cap W$, it thus suffices to show $N_{\mathit{reg}}$ is nonempty.

Since $f$ is not identically zero, not all derivatives (of arbitrary order) of $f$ vanish identically on $N$. If some first order derivative of $f$ does not vanish identically on $N$, let $h=f$. Otherwise, suppose $k$ is such that a derivative of $f$ of order $k$ does not vanish identically on $N$, and all derivatives of $f$ of order less than $k$ vanish identically on $N$. Let $h$ be one of the derivatives of order $k-1$. We obtain a function $h\colon U\to\C$, holomorphic, vanishing on $N$, and such that without loss of generality the $z_n$ derivative does not vanish identically on $N$. Then there is some point $p\in N$ such that $\frac{\partial h}{\partial z_n}(p)\ne0$. We apply the implicit function theorem at $p$ to find $g$ such that
\[h\bigl(z_1,\ldots,z_{n-1},g(z_1,\ldots,z_{n-1})\bigr)=0,\]
and $z_n=g(z_1,\ldots,z_{n-1})$ is the unique solution to $h=0$ near $p$.

The zero set of $h$ contains $N$, the zero set of $f$. We must show equality near $p$. That is, we need to show that near $p$, every zero of $h$ is also a zero of $f$. Write $p=(p',p_n)$. The function $\xi\mapsto f(p',\xi)$ has an isolated zero in a small disc $\Delta$ around $p_n$ and is nonzero on the circle $\partial\Delta$. By Rouch\'e's theorem, $\xi\mapsto f(z',\xi)$ has a zero in $\Delta$ for all $z'$ sufficiently close to $p'$ (close enough to make $\sabs{f(p',\xi)-f(z',\xi)}<\sabs{f(p',\xi)}$ for all $\xi\in\partial\Delta$). Since $g(z')$ is the unique solution $z_n$ to $h(z',z_n)=0$ near $p$ and the zero set of $f$ is contained in the zero set of $h$, we are done.
\end{proof}
\begin{theorem}[Injective holomorphic mappings]
Suppose $U\subset\C^n$ is an open set and $f\colon U\to\C^n$ is holomorphic and one-to-one. Then the Jacobian determinant is never equal to zero on $U$.
In particular, if a holomorphic map $f\colon U\to V$ is one-to-one and onto for two open sets $U,V\subset\C^n$, then $f$ is biholomorphic.
\end{theorem}
The function $f$ is locally biholomorphic, in particular $f^{-1}$ is holomorphic, on the set where the Jacobian determinant
\[J_f(z)=\det Df(z)=\det\left[\frac{\partial f_k}{\partial z_\ell}(z)\right]_{k\ell}\]
is not zero. This follows from the inverse function theorem, which is just a special case of the implicit function theorem. The trick to prove the theorem above is to prove that $J_f$ is nowhere zero.

In one complex dimension, every nonconstant holomorphic function $f$ can, in the proper local holomorphic coordinates (and up to adding a constant), be written as $z^d$ for $d=1,2,\ldots$: Near a $z_0\in\C$, there exists a constant $c$ and a local biholomorphic $g$ with $g(z_0)=0$ such that $f(z)=c+\bigl(g(z)\bigr)^d$. So $f$ is one-to-one precisely if $d=1$.
\begin{proof}[Proof of the theorem]
We proceed by induction. We know the theorem for $n=1$. Suppose $n>1$ and suppose we know the theorem is true for dimension $n-1$.

Suppose for contradiction that $J_f=0$ somewhere. First suppose that $J_f$ is not identically zero. Find a regular point $q$ of the zero set of $J_f$. Write the zero set of $J_f$ near $q$ as
\[z_n=g(z_1,\ldots,z_{n-1})\]
for some holomorphic $g$. If we prove the theorem near $q$, we are done. Without loss of generality assume $q=0$. The biholomorphic (near the origin) map
\[\Psi(z_1,\ldots,z_n)=\bigl(z_1,z_2,\ldots,z_{n-1},z_n-g(z_1,\ldots,z_{n-1})\bigr)\]
takes the zero set of $J_f$ to the set given by $z_n=0$. Considering $f\circ\Psi^{-1}$ instead of $f$, we assume that $J_f=0$ on the set given by $z_n=0$. We also assume that $f(0)=0$.

If $J_f$ vanishes identically, then there is no need to do anything other than a translation. In either case, we may assume that $0\in U$, $f(0)=0$, and $J_f=0$ when $z_n=0$. Really, all we need is for the set where $J_f=0$ to be a sufficiently large set.

We wish to show that all the derivatives of $f$ in the $z_1,\ldots,z_{n-1}$ variables vanish whenever $z_n=0$. This would clearly contradict $f$ being one-to-one, as $f(z_1,\ldots,z_{n-1},0)$ would be constant. So for any point on $z_n=0$, consider one of the components of $f$ and one of the derivatives of that component. Without loss of generality, suppose the point is $0$, and for contradiction suppose $\frac{\partial f_1}{\partial z_1}(0)\ne0$. The map
\[G(z_1,\ldots,z_n)=\bigl(f_1(z),z_2,\ldots,z_n\bigr)\]
is biholomorphic on a small neighborhood of the origin. The function $f\circ G^{-1}$ is holomorphic and one-to-one on a small neighborhood. By the definition of $G$,
\[f\circ G^{-1}(w_1,\ldots,w_n)=\bigl(w_1,h(w)\bigr),\]
where $h$ is a holomorphic mapping taking a neighborhood of the origin in $\C^n$ to $\C^{n-1}$. The mapping
\[\varphi(w_2,\ldots,w_n)=h(0,w_2,\ldots,w_n)\]
is a one-to-one holomorphic mapping of a neighborhood of the origin in $\C^{n-1}$ to $\C^{n-1}$. By the induction hypothesis, the Jacobian determinant of $\varphi$ is nowhere zero.

If we differentiate $f\circ G^{-1}$, we notice $D(f\circ G^{-1})=Df\circ D(G^{-1})$. So at the origin
\[\det D(f\circ G^{-1})=(\det Df)(\det D(G^{-1}))=0.\]
We obtain a contradiction, as at the origin
\[\det D(f\circ G^{-1})=\det D\varphi\ne0.\qedhere\]
\end{proof}
\subsection*{整理说明（非作者正文）}
保留作者的正则点构造、Rouch\'e 步骤、临界集坐标展平及维数归纳；无关例子、习题及后续 Jacobian 猜想讨论不收。标准先修为一复变局部零点正规形、全纯隐/逆函数定理、Rouch\'e 定理、恒等定理和链式法则。难度在于把可能奇异的 Jacobian 零集转换为可展平的局部超曲面，再把单射性限制到低维映射得到矛盾；不能从单射性直接引用待证非退化结论。本题不同于033的自同构线性化、058的适当映射不存在性。原 $D(f\circ G^{-1})$ 的简记按正文保留。数学 TeX 按原生正文摘取并作断行，不声称取得整书原件。
\subsection*{可修改的简短标注（非作者正文）}
$c$：全纯单射映射的局部可逆性

$a$：解析零集正则点；最低非消失导数；Rouché定理；局部坐标展平；隐函数定理；维数归纳

\texttt{cross\_theory\_status = yes}。将 Jacobian 退化问题转化为解析零集的几何结构，以局部坐标把该集展平，再用低维单射映射的非退化性排除临界超曲面。位置：所附正则点定理与主定理归纳证明。
标签为非穷尽、可修改初稿。
\subsection*{许可与修改记录}
原作者及作品见来源栏；来源采用 CC BY-SA 4.0。本题证摘录和附加整理沿用该许可。2026-09-28：增加题号、来源定位、独立排版、整理说明和初稿标注；数学内容的取材范围及排版变化见整理说明。
\end{document}
